With diagnostic labels randomly reassigned across participants (5 permutations; each participant keeps one label, but which one is random), no disorder signal remains and an honest estimate must sit at 0.5 (Table~\ref{tab:e7}). On DEPRESJON, subject-wise AUROC on permuted labels was 0.508 (logistic regression) and 0.517 (XGBoost), whereas record-wise AUROC was 0.646 (range over permutations 0.616 to 0.689) and 0.744 (0.721 to 0.768), reaching 0.740 and 0.898 at participant level. On PSYKOSE, subject-wise AUROC on permuted labels was 0.525 (logistic regression) and 0.506 (XGBoost), whereas record-wise AUROC was 0.638 (range over permutations 0.554 to 0.714) and 0.733 (0.701 to 0.799), reaching 0.719 and 0.895 at participant level. On HYPERAKTIV, subject-wise AUROC on permuted labels was 0.528 (logistic regression) and 0.533 (XGBoost), whereas record-wise AUROC was 0.621 (range over permutations 0.591 to 0.633) and 0.671 (0.641 to 0.703), reaching 0.666 and 0.753 at participant level. Under record-wise splitting the models therefore achieved discrimination comparable to the \emph{genuine} subject-wise results of E1 on labels that carry no information about the disorder at all: the design measures the model's ability to recognise a participant from other days of the same participant, and nothing else. The effect was larger for XGBoost than for logistic regression in every cohort, consistent with the E1 ordering.
